\documentclass[10pt]{article}
\usepackage[a4paper,margin=1in]{geometry}
\usepackage{amsmath,amssymb,hyperref}
\title{A counterexample to Conjecture 00000004396}
\author{gaochengzhecpu}
\date{October 3, 2026}
\setlength{\emergencystretch}{1em}
\hypersetup{hidelinks}
\begin{document}
\maketitle

The conjecture asserts, in particular, that every M\"obius ladder is
Sidorenko. This assertion is false for the eight-vertex M\"obius ladder.
No restriction to bipartite ladders or to a congruence class of their orders
appears in either language of the source.

For finite simple graphs $H,G$, write
\[
 t(H,G)=\frac{\operatorname{hom}(H,G)}{|V(G)|^{|V(H)|}}.
\]
Here a homomorphism is an arbitrary vertex map preserving every edge;
it need not be injective. The Sidorenko property is
\[
 t(H,G)\geq t(K_2,G)^{|E(H)|}\qquad\text{for every nonempty finite graph }G.
\]
This is the standard finite-graph formulation\cite{sidorenko}.

Let $H$ have vertex set $\mathbb Z/8\mathbb Z$, with edges between
consecutive vertices and between opposite vertices. Thus $H$ is the cycle
$C_8$ together with the four edges $\{i,i+4\}$, $0\leq i<4$.
It has twelve edges and is a M\"obius ladder independently of conventions
about including the smallest four-vertex member of this family.

The vertices $0,1,2,3,4,0$ form a cycle of length five in $H$.
A homomorphism $H\to K_2$ would alternate its two values along every edge
of this cycle, which is impossible. Consequently
\[
 \operatorname{hom}(H,K_2)=0,\qquad t(H,K_2)=0.
\]
On the other hand, $K_2$ has two ordered adjacent pairs out of four
ordered vertex pairs, so
\[
 t(K_2,K_2)=\frac12,\qquad
 t(K_2,K_2)^{|E(H)|}=2^{-12}>0.
\]
The required inequality is therefore false. This refutes the conjunction
in the source; no claim about odd subdivisions of trees is needed.

\paragraph{Formal verification.}
The self-contained Lean 4.19.0 file \texttt{lean/Main.lean} defines actual
finite simple graphs, edge-preserving vertex maps, and the Sidorenko
inequality. Counts use complete duplicate-free lists of maps. The
inequality is written after multiplication by positive integer
denominators, so there is no approximation or division convention:
\[
 (2|E(G)|)^{|E(H)|}|V(G)|^{|V(H)|}
 \leq \operatorname{hom}(H,G)|V(G)|^{2|E(H)|}.
\]
The proof first excludes every map $\operatorname{Fin}8\to
\operatorname{Fin}2$ directly using the five-cycle. It then proves that
the 256 binary codes enumerate all such functions, each exactly once,
using inverse encoding and decoding maps. The computed counts are thus
counts of the intended objects. The final theorem
\texttt{conjecture\_\allowbreak{}4396\_\allowbreak{}false} negates the universal assertion about
the cycle-plus-opposite-edge family. The file uses no \texttt{sorry},
\texttt{admit}, \texttt{native\_\allowbreak{}decide}, or additional axiom.

\paragraph{Previous withdrawn submission.}
An earlier, independently authored submission by \texttt{orionsheep}
used the smaller $K_4$ member and the same necessary bipartiteness
observation. It was withdrawn by its author, with no review rejection
recorded, in pull request
\href{https://github.com/The-Last-Math-Competition/The-Last-Math-Competition/pull/318}{\#318}.
The present package supplies an independent eight-vertex construction
and full finite-model verification; it does not claim priority for the
general bipartiteness obstruction.

\begin{thebibliography}{1}
\bibitem{sidorenko}
\emph{Sidorenko's conjecture, colorings and independent sets},
\url{https://arxiv.org/abs/1603.05888}.
\end{thebibliography}
\end{document}
